\documentclass[10pt,a4paper]{amsart}
\usepackage[margin=19mm]{geometry}
\usepackage{amsmath,amssymb,mathtools}
\usepackage[scr=rsfs,cal=euler]{mathalfa}
\usepackage{xfrac}
\usepackage[unicode,hidelinks,bookmarks=false]{hyperref}
\newcommand{\ie}{i.e.\ }
\newcommand{\cf}{cf.\ }
\newcommand{\resp}{respectively\ }
\newcommand{\Subsection}{\subsection}
\providecommand{\nobreakdash}{\nobreak\hbox{-}}
\setlength{\emergencystretch}{2em}
\multlinegap0pt
\usepackage[matrix,arrow,tips,curve]{xy}
\usepackage{enumerate}
\providecommand{\noop}[1]{}
\providecommand{\bysame}{\leavevmode\hbox to3em{\hrulefill}\thinspace}
\renewcommand{\qedsymbol}{$\blacksquare$}
\newtheorem{thm}{Theorem}[section]
\newtheorem{lemma}[thm]{Lemma}
\newtheorem{thmdef}[thm]{Theorem - Definition}
\newtheorem{corollary}[thm]{Corollary}
\newtheorem{proposition}[thm]{Proposition}
\newtheorem*{proposition817}{Proposition 8.17}
\newtheorem{question}[thm]{Question}
\newtheorem{thmdefi}[thm]{Theorem - Definition}
\newtheorem{step}{Step}
\newtheorem{prgdim}[thm]{}
\newtheorem*{thm*}{Theorem}
\theoremstyle{definition}
\newtheorem{definition}[thm]{Definition}
\newtheorem{setup}[thm]{Set-up}
\newtheorem{property}[thm]{}
\newtheorem{example}[thm]{Example}
\newtheorem{remark}[thm]{Remark}
\newtheorem{remdefi}[thm]{Definition - Remark}
\newtheorem{prg}[thm]{}
\newtheorem{parg}{}[thm]
\newtheorem{subparg}{}[parg]
\newtheorem{notation}[thm]{Notation}
\renewcommand{\theequation}{\thethm}
\numberwithin{figure}{section}
\numberwithin{table}{section}
\newcommand{\ph}{\varphi}
\newcommand{\w}{\widetilde}
\newcommand{\ma}{\mathcal}
\newcommand{\la}{\longrightarrow}
\newcommand{\ol}{\mathcal{O}}
\newcommand{\wi}{\widehat}
\newcommand{\pr}{\mathbb{P}}
\newcommand{\Q}{\mathbb{Q}}
\newcommand{\C}{\mathbb{C}}
\newcommand{\R}{\mathbb{R}}
\newcommand{\Z}{\mathbb{Z}}
\newcommand{\N}{\mathcal{N}_1}
\newcommand{\Nu}{\mathcal{N}^1}
\newcommand{\Gr}{\operatorname{Gr}}
\newcommand{\dom}{\operatorname{dom}}
\newcommand{\depth}{\operatorname{depth}}
\newcommand{\Spec}{\operatorname{Spec}}
\newcommand{\Sing}{\operatorname{Sing}}
\newcommand{\Pic}{\operatorname{Pic}}
\newcommand{\SQM}{\operatorname{SQM}}
\newcommand{\Rat}{\operatorname{Ratcurves}^n}
\newcommand{\Cox}{\operatorname{Cox}}
\newcommand{\im}{\operatorname{Im}}
\newcommand{\NE}{\operatorname{NE}}
\newcommand{\Exc}{\operatorname{Exc}}
\newcommand{\Bs}{\operatorname{Bs}}
\newcommand{\Supp}{\operatorname{Supp}}
\newcommand{\Lo}{\operatorname{Locus}}
\newcommand{\codim}{\operatorname{codim}}
\newcommand{\Eff}{\operatorname{Eff}}
\newcommand{\Nef}{\operatorname{Nef}}
\newcommand{\Chow}{\operatorname{Chow}}
\newcommand{\Mov}{\operatorname{Mov}}
\newcommand{\MCD}{\operatorname{MCD}}
\newcommand{\mov}{\operatorname{mov}}
\newcommand{\Hom}{\operatorname{Hom}}
\newcommand{\Homb}{\operatorname{Hom}_{\text{bir}}}
\newcommand{\Hilb}{\operatorname{Hilb}}
\newcommand{\Aut}{\operatorname{Aut}}
\newcommand{\Bl}{\operatorname{Bl}}
\newcommand{\Id}{\operatorname{Id}}
\newcommand{\sm}{\text{\it sm}}
\newcommand{\pt}{\text{\it pt}}
\newcommand{\pts}{\text{\it pts}}
\newcommand{\reg}{\text{\it reg}}
\newcommand{\s}{\scriptscriptstyle}
\title[Fano $4$-folds with $b_2\geq 7$]{Towards the classification of Fano $4$-folds with $b_2\geq 7$: the case of every fixed prime divisor of type $(3,0)^{\sm}$ (Proposition 9.4)}
\author{C.~Casagrande}
\date{Source: arXiv:2508.21207v1 (CC BY 4.0)}
\begin{document}
\begin{abstract}
  We study (smooth, complex) Fano $4$-folds $X$ with Picard number $\rho_X\geq 7$. We show that if $\rho_X>9$, then $X$ is a product of del Pezzo surfaces (Th.~1.1), thus improving results in [Cas24, CS24]; the statement is now optimal.  In the range $\rho_X\in\{7,8,9\}$ we show that if $X$ is not a product of surfaces, and has no small elementary contraction, then it is the blow-up of a cubic $4$-fold along a special configuration of planes
  (Th.~1.2). When instead $\rho_X\geq 7$ and $X$ has a small elementary contraction, we study $X$ depending on its fixed prime divisors, giving explicit results on the geometry of $X$ in the framework of birational geometry.
  In particular for the boundary case $\rho_X=9$ we show that either $X$ is a product of surfaces, or $X$ belongs to two explicit families, or there is a sequence of flips $X\dasharrow X'$ where $X'$ is a smooth projective $4$-fold with an elementary contraction onto a $3$-fold (Th.~1.4).

  In the paper we also give several results on rational contractions of fiber type of Fano $4$-folds, and more generally of Mori dream spaces; in particular we use some properties of del Pezzo surfaces over non-closed fields, applied to generic fibers.
\end{abstract}
\maketitle
\noindent\emph{Source and licence.} Towards the classification of Fano $4$-folds with $b_2\geq 7$. Version arXiv:2508.21207v1 (CC BY 4.0), distributed under the Creative Commons Attribution 4.0 International licence, http://creativecommons.org/licenses/by/4.0/, as recorded in the arXiv OAI licence metadata for this exact version and shown on the abstract page https://arxiv.org/abs/2508.21207v1. Full rights notice: licenses/arxiv-2508.21207-CC-BY-4.0.txt.\par\smallskip
\noindent\textit{Editorial scope.} Counted unit: the proposition that carries the whole case in which every fixed prime divisor is of type $(3,0)^{\sm}$, source label all30 (source0.tex lines 2410--2412), printed as Proposition 9.4 in the published version and as Proposition 9.4 here, delivered together with its complete original author proof at source lines 2413--2606 (194 lines) as the single primary proof body. The proof is an adjacent block: it opens with the paper's own \textbackslash{}begin\{proof\} on the line after the statement, closes with \textbackslash{}end\{proof\} at source line 2606, ends with the author's own \textbackslash{}qedhere and its last sentence names its own statement (This concludes the proof of Prop.~all30). The paper already contains a paragraph that announces this proposition as Prop.~9.4 (source line 2391). No mathematics is written, repaired or re-derived: the statement and the proof are the paper's own text line for line, including its own xymatrix diagrams, its own stepcounter steps, its own numbered environments and its own printed slips. The paper is already an amsart article (source0.tex line 2, class options a4paper and 12pt), so no publisher class is replaced or shipped: the wrapper is the lane's pinned amsart editorial wrapper, and the paper's own settings are copied verbatim from the source --- its \textbackslash{}newtheorem declarations (source lines 26--48, that is the section-numbered thm counter shared by lemma, thmdef, corollary, proposition, question, thmdefi, prgdim, definition, setup, property, example, remark, remdefi, prg, parg, subparg and notation, the starred proposition817 and thm*, the step counter, the definition theorem style, and the empty names of the prg and property environments), its equation renumbering \textbackslash{}renewcommand\{\textbackslash{}theequation\}\{\textbackslash{}thethm\} (source line 50), its \textbackslash{}numberwithin settings for figure and table (source lines 53--54), its proof-end symbol (source line 20), its macro block (source lines 56--101) and its \textbackslash{}bysame and \textbackslash{}noop definitions (source lines 3273--3274), and the xy and enumerate packages are loaded as the paper loads them. The paper's own \textbackslash{}usepackage\{times\} load is not reproduced: it selects a legacy eight-bit Times family that the pinned runtime does not provide for xelatex, so the wrapper keeps its own Latin Modern text font; this is a typographic substitution of the wrapper that changes no mathematics. Consequently every printed number of the delivered unit is produced by the paper's own counters and coincides with the published version. Because the eight sections that precede it are not part of this unit, the wrapper initializes the section counter editorially with \textbackslash{}setcounter\{section\}\{8\} immediately before the delivered span, so that the paper's own section heading prints as Section 9 and the whole numbering of the unit coincides with the published one: the remark united prints 9.1, the paper's own main theorem on this case prints Theorem 9.2, the overview paragraph prints 9.3, the counted statement prints Proposition 9.4, and inside the counted proof the author's own counter steps and numbered environments print 9.5 (the numbered equation labelled sections), 9.6 (rhoY1), 9.7 (dimZ2), 9.8 (avocado), 9.10 (flixbus), 9.11 (duo), 9.13 (colorado), 9.14, 9.15, 9.17 (verylast), 9.18 (sections2) and 9.19 (summer), with the author's own unlabelled counter steps carrying the remaining numbers; the published version itself prints the same numbers (for instance see 9.7 - 9.17 at source line 2403 and by Th. 7.1 at source line 2408). Required context retained uncounted: the whole of Section 9 up to the counted statement (source lines 2372--2409), that is the paper's own section heading for this case, the remark united showing that a Fano 4-fold with rho at least 7 and a small elementary contraction must have a fixed prime divisor of type (3,0)sm, (3,1)sm or (3,0)Q, the paper's own main theorem on this case with its two alternatives, and the overview paragraph of its proof, which is also where the counted proposition is announced and which refers to the counted proposition, to its companion proposition on the other case, to the paper's earlier theorem on elementary rational contractions onto 3-folds and to its section on Fano models. Every definition, theorem, remark, numbered display, footnote and citation of those source lines is retained. Note that the proof of Theorem 9.2 itself is not part of this unit: it is a short detached block at source lines 2651--2655, after the companion proposition, and it is not delivered, so Theorem 9.2 appears here as an announced statement with the counted proposition supplying the case it describes. The paper's abstract (source lines 137--141) is reproduced in the wrapper front matter with its three cross-references and its two out-of-unit citations printed as their published printed numbers (Th. 1.1, Th. 1.2, Th. 1.4 and [Cas24, CS24]), because front matter is not part of the source-span machinery; nothing else of the abstract is changed. Not part of this unit and not redistributed: the paper's address and email front matter (source lines 118--122); Sections 1--8 (source lines 144--2371) with the Introduction, the notation and preliminaries, and all the preparatory sections on rational contractions, on factoring rational contractions through 3-folds, on fixed prime divisors, on quasi-elementary contractions onto surfaces and on fixed divisors of type (3,2); the rest of the paper from the companion proposition onwards (source lines 2607--3462), that is the case in which there are also fixed prime divisors of other types, the short proof of Theorem 9.2, the closing sections on divisors of type (3,1)sm or (3,0)Q and the examples; and the paper's bibliography listing, of which only the cited entries are transcribed. The delivered text cites four keys (blowup, fanoEMS, kollar and GraberHarrisStarr, thirteen citation commands in all), and those four entries are transcribed verbatim from the paper's own in-source thebibliography (source lines 3281--3456) with the paper's own printed bibliography labels preserved (Cas17, IP99, Kol96 and GHS03, exactly as the published list prints them), so every printed citation of the delivered unit is the published one, including the citation of the paper's own earlier work in the optional argument of the counted statement (refinement of [Cas17], Th. 4.3). The entry Cas17 is transcribed exactly as the paper's bibliography prints it, that is beginning with the paper's own \textbackslash{}bysame rule, the typographic stand-in for the author name repeated from the preceding Casagrande entries; because that preceding entry is not part of this unit, the delivered list shows the rule at the head of the entry, which is the paper's own printed form and is kept literally rather than expanded. Reference adaptation: fourteen keys, twenty-one occurrences, method reference\_only\_adaptation recorded in delivery-evidence.json. Nine of them lie in the retained context (summary prints 4.10, kawamata 2.13, dim32 5.3, fixed 5.1, 30intro 1.3, CS 1.5, one30 9.20, fanomodel 11.1 and sharp32 7.1) and eight in the counted proof (CS 1.5, fixedprime 2.3, h0 2.12, fanomodel 11.1, specialsurf 3.19, calanques 3.12, kawamata 2.13 and malpensa 8.34); each occurrence is delivered as a \textbackslash{}href link to the abstract page of this exact version with the printed number of the target retained, so the delivered unit compiles with no unresolved reference, and every printed number was checked against the published PDF. All other labels and references of the retained spans are the paper's own and resolve inside the unit (sec30, united, 30, overview30, all30, sections, rhoY1, dimZ2, avocado, flixbus, duo, colorado, Y1, Y2, verylast, sections2 and summer), and no label occurs twice. Figures: the paper has six figure environments whose bodies are \textbackslash{}input\{figura*.pdftex\_t\} picture files (source lines 1175--1182, 1273--1278, 2921--2926, 2991--2996, 3083--3088 and 3124--3129). None of them lies inside the retained spans and nothing in the retained spans refers to them, so no unresolved \textbackslash{}input and no undeclared figure dependency reaches the compile, and the six pictures of the paper are not redistributed. The paper's loads of tikz-cd, graphicx, babel, mathrsfs and bm are unused by the retained spans and are not reproduced; its inputenc with latin1 is dropped because xelatex rejects inputenc and the whole source is pure ASCII (the paper has no non-ASCII character, so no character had to be replaced, dropped or re-encoded and no fidelity receipt for a character is needed), which the xelatex receipt confirms with an empty missing-character list. The paper's own page-layout settings (\textbackslash{}setlength\{\textbackslash{}textwidth\}, \textbackslash{}calclayout and the section and subsection patches at source lines 103--106 and 129--130) are replaced by the wrapper's pinned geometry, and the wrapper's pinned mathalfa options set the calligraphic alphabet in Euler script where the paper used the default one; both are page-layout and typographic decisions of the wrapper that change no mathematics. Printed slips kept literal: no printed word, symbol, hypothesis or number of the counted statement or of its proof was altered; the proof keeps the paper's own spelling, its own use of ibid. inside brackets, its own spacing commands such as mskip and medskip, and its own notation for symmetric and asymmetric powers. No mathematics is written, repaired or re-derived; all mathematics of the unit is native text with no image dependency.
\setcounter{section}{8}
\section{Fano $4$-folds with a  divisor of type $(3,0)^{\sm}$}\label{sec30}
\noindent Let $X$ be a Fano $4$-fold with $\rho_X\geq 7$, that is not a product of surfaces; we consider now the case where $X$ has some small elementary contraction.  The following remark shows that such $X$ must contain a fixed prime divisor  of type $(3,0)^{\sm}$, $(3,1)^{\sm}$, or $(3,0)^Q$.
 \begin{remark}\label{united}
Let $X$ be a smooth Fano $4$-fold with $\rho_X\geq 7$, having a small elementary contraction $f\colon X\to Y$. We have $\NE(f)\not\subset\mov(X)=\Eff(X)^{\vee}$, therefore there exists a one-dimensional face $\tau$ of $\Eff(X)$ such that $\tau\cdot\NE(f)<0$. By Cor.~\href{https://arxiv.org/abs/2508.21207v1}{4.10} $\tau$ is a fixed face, hence $\tau=\R_{\geq 0}[D]$ for some  fixed prime divisor $D\subset X$, and $D\cdot\NE(f)<0$. In particular $D$ contains $\Lo(f)$ which is a union of exceptional planes (Lemma \href{https://arxiv.org/abs/2508.21207v1}{2.13}), and $D$ cannot be of type $(3,2)$ (Lemma \href{https://arxiv.org/abs/2508.21207v1}{5.3}). We conclude that $D$ must be of type $(3,0)^{\sm}$, $(3,1)^{\sm}$, or $(3,0)^Q$ (see Th.-Def.~\href{https://arxiv.org/abs/2508.21207v1}{5.1}).
          \end{remark}

In this section we consider the case where $X$ has a fixed prime divisor of type $(3,0)^{\sm}$, and show the following result (Th.~\href{https://arxiv.org/abs/2508.21207v1}{1.3} from the Introduction).
\begin{thm}\label{30}
Let $X$ be a smooth Fano $4$-fold with $\rho_X\geq 7$ having a fixed prime divisor of type $(3,0)^{\sm}$. Then one of the following holds:
  \begin{enumerate}[$(i)$]
  \item  $X$ has an elementary rational contraction onto a $3$-fold and $\rho_X\leq 9$;
  \item $\rho_X=7$,
$h^0(X,-K_X)\in\{2,\dotsc,15\}$,
and there is a SQM $X\dasharrow X'$ such that $X'=\Bl_{6\mskip1mu\pts}Y$, $Y$ a smooth Fano $4$-fold with $\rho_Y=1$ and $h^0(Y,-K_Y)=h^0(X,-K_X)+90\in\{92,\dotsc, 105\}$. Moreover every fixed prime divisor of $X$ is of type $(3,0)^{\sm}$.
    \end{enumerate}
\end{thm}
\begin{prg}[\em Overview of the proof of Th.~\ref{30}]\label{overview30}
The proof follows the same strategy as in \cite{blowup}, where the bound $\rho_X\leq 12$ is shown. Let us note that, in case $(i)$ where $X$ has an elementary rational contraction onto a $3$-fold, the bound $\rho_X\leq 9$ follows from Th.~\href{https://arxiv.org/abs/2508.21207v1}{1.5}.

We treat separately the two cases where every fixed prime divisor is of type $(3,0)^{\sm}$ (Prop.~\ref{all30}), and where there are also fixed prime divisors of other types (Prop.~\href{https://arxiv.org/abs/2508.21207v1}{9.20}).

In the first case, there are a smooth Fano $4$-fold $Y$ and a SQM $X\dasharrow X'$ such that $X'=\Bl_{\pts}Y$, and $Y$ does not have divisorial elementary rational contractions. We compute that $h^0(Y,-K_Y)=h^0(X,-K_X)+15(\rho_X-\rho_Y)$.  

If $Y$ has an elementary rational contraction onto a $3$-fold, then we lift this contraction to $X$, and get $(i)$.
If $\rho_Y=1$, then  either $Y\cong\pr^4$ and we get again $(i)$ (see \S\href{https://arxiv.org/abs/2508.21207v1}{11.1}), or we get $h^0(Y,-K_Y)\leq 105$ from \cite{blowup}, which implies
$(ii)$ (see \ref{rhoY1}).

Otherwise, we show that we can reduce to the situation where
 $\rho_Y=2$ and $Y$ has two distinct elementary rational contractions of fiber type $Y\dasharrow B_i$ where $B_i$ is $\pr^1$ or $\pr^2$, for $i=1,2$; moreover $\rho_X-\rho_Y\geq 5$ gives 
$h^0(Y,-K_Y)\geq 75$.
We show that this case does not happen, by studying explicitly the geometry of $Y$.
 This is the longest part of the proof (see \ref{dimZ2} - \ref{verylast}).

\medskip

Consider now the case where $X$ has a fixed prime divisor of type $(3,0)^{\sm}$ and also some fixed prime divisor of another type. We know that $\Eff(X)$ is generated by classes of fixed prime divisors, and that every $2$-dimensional face is fixed (Cor.~\href{https://arxiv.org/abs/2508.21207v1}{4.10}). Then there must a  fixed prime divisor $D$ of type $(3,0)^{\sm}$ and a fixed prime divisor $B$, not of type $(3,0)^{\sm}$, that are adjacent. We consider the contraction $X\dasharrow Y$ of $D$, so that $Y$ is a smooth Fano $4$-fold, and the transform $B_Y\subset Y$ is still a fixed prime divisor. By analysing 
 the possibilities for $B_Y$, we show that in each case there is a fixed prime divisor $E\subset X$, of type $(3,2)$, such that $D\cap E=\emptyset$. Then we get $(i)$ by Th.~\href{https://arxiv.org/abs/2508.21207v1}{7.1}.
\end{prg}

\begin{proposition}[refinement of \cite{blowup}, Th.~4.3]\label{all30}
  Let $X$ be a smooth Fano $4$-fold with $\rho_X\geq 7$, and suppose that every fixed prime divisor of $X$ is of type $(3,0)^{\sm}$. Then the statement of Th.~\ref{30} holds.
  \end{proposition}

  \begin{proof}
    Let us first note that, if $X$ has an elementary rational contraction onto a $3$-fold, then $\rho_X\leq 9$ by Th.~\href{https://arxiv.org/abs/2508.21207v1}{1.5}, and we get $(i)$.

    We follow  \cite[proof of Lemma 4.22]{blowup}, where it is shown that there is a diagram
    $$X\stackrel{\ph}{\dasharrow} Y\stackrel{f}{\la} Z$$
    where $\ph$ is a sequence of flips and divisorial elementary  contractions, and $f$ is an elementary contraction of fiber type.

    Up to increasing the number of flips and divisorial contractions in $\ph$, we can also assume that $Y$ has no divisorial elementary  rational contractions, hence $Y$ contains no fixed prime divisors (see Rem.~\href{https://arxiv.org/abs/2508.21207v1}{2.3}). 

    By [\emph{ibid.}, Lemma 4.19], up to SQM we can assume that   $Y$ is a smooth Fano $4$-fold, and  there is a resolution
    $\ph'\colon X'\to Y$
    of $\ph$ such that $\ph'$ 
 is  the blow-up of $Y$ in $r:=\rho_X-\rho_Y$ distinct points. Then we have:
    \stepcounter{thm}
    \begin{equation}\label{sections}
     2\leq h^0(X,-K_X)=h^0(Y,-K_Y)-15r,
   \end{equation}
   by Th.~\href{https://arxiv.org/abs/2508.21207v1}{2.12} and [\emph{ibid.}, Prop.~3.3].

   If $\dim Z=3$, then by [\emph{ibid.}, Lemma 4.21] we have $(i)$.
    \begin{prg}\label{rhoY1}
      Suppose that $\rho_Y=1$ and $Z=\{\pt\}$.
      If $Y\cong\pr^4$,
      then  $X$ is the Fano model of $\Bl_{\pts}\pr^4$ and has an elementary rational contraction onto $\Bl_{\pts}\pr^3$ (see \S\href{https://arxiv.org/abs/2508.21207v1}{11.1}), thus we have $(i)$.
      If instead $Y\not\cong\pr^4$, 
 as in [\emph{ibid.}, proof of Th.~3.2, p.~380] we see that $h^0(Y,-K_Y)\leq 105$, which together with \eqref{sections} implies $(ii)$.
\end{prg}
  \begin{prg}\label{dimZ2}
   Suppose that $\dim Z=2$, and note that $Z$ is smooth (Lemma \href{https://arxiv.org/abs/2508.21207v1}{3.19}).   
    We notice that
    $Z$ cannot have divisorial  elementary contractions, otherwise the pullback under $f$ of the exceptional divisor would be a fixed prime divisor in $Y$, against our reductions.
Therefore $Z$ is a minimal smooth rational surface  with no 
divisorial elementary contraction, and we conclude that either $Z\cong\pr^2$ and $\rho_Y=2$, or $Z\cong\pr^1\times\pr^1$ and $\rho_Y=3$. We also note that if $\dim Z=1$, then $Z\cong\pr^1$ and $\rho_Y=2$.
\end{prg}
  \begin{prg}\label{avocado}
   Suppose that $Z\cong\pr^1\times\pr^1$, and consider the blow-up
of the first point $W:=\Bl_{p_1}Y\to Y$ in $\ph'\colon X'\to Y$. The composition $W\to\pr^1\times\pr^1$ is not equidimensional, and by Prop.~\href{https://arxiv.org/abs/2508.21207v1}{3.12} it factors as
    $$\xymatrix{{W}\ar[r]\ar@{-->}[d]&Y\ar[d]\\
     S\ar[r]&{\pr^1\times\pr^1}
    }$$
    where $W\dasharrow S$ is an elementary rational contraction, $S$ is a smooth surface (see Lemma \href{https://arxiv.org/abs/2508.21207v1}{3.19}), and $S\to\pr^1\times\pr^1$ is the blow-up of a smooth point, thus $S\cong\Bl_{2\mskip1mu\pts}\pr^2$.

    Consider a $(-1)$-curve of $S$ contracted by $S \to \pr^2$, and its pullback $G$ in $W$. This is a fixed
prime divisor, and up to to replacing  $W$ with a SQM, we can assume that $G = \Exc(g)$ for some
divisorial elementary contraction $g \colon W \to Y'$ (see Rem.~\href{https://arxiv.org/abs/2508.21207v1}{2.3}). Then $g$ must be the blow-up of a smooth point (by \cite[Lemma 4.19(2)]{blowup} applied to the composition $X\dasharrow Y'$), and there is an elementary rational contraction $Y'\dasharrow \mathbb{F}_1$.
$$\xymatrix{{W}\ar[r]^g\ar@{-->}[d]&{Y'}\ar@{-->}[d]\\
S\ar[r]&{\mathbb{F}_1}
}$$
Now we consider the blow-up $\mathbb{F}_1\to\pr^2$ and the composition $Y'\dasharrow \pr^2$; it factors again as an elementary divisorial rational contraction $Y'\dasharrow Y''$ and an elementary contraction $Y''\to\pr^2$. Finally, by replacing $Y$ with $Y''$,
 and we can assume that $\rho_Y=2$ and $Z\cong\pr^2$.
\end{prg}
  \begin{prg}
    We are left with the possibility that $\rho_Y=2$ and $Y$ has two distinct elementary rational contractions of fiber type $Y\dasharrow B_i$  where $B_i$ is $\pr^1$ or $\pr^2$, for $i=1,2$, so that $\Mov(Y)=\Eff(Y)$ and the two one-dimensional faces of this cone are the pullbacks of $\Nef(B_i)$, $i=1,2$.
    We are going to show that
   this contradicts $\rho_X\geq 7$.
   \end{prg}
 \begin{prg}\label{flixbus}
Suppose that $Y$ has index $>1$. Then $Y$ cannot have small elementary contractions (see Lemma \href{https://arxiv.org/abs/2508.21207v1}{2.13}), therefore 
both maps $Y\to B_i$ must be regular. Since each map must be finite on the fibers of the other one, such fibers are at most $2$-dimensional, and we conclude that $Y$ has two elementary contractions onto $\pr^2$.

Fano $4$-folds with index $>1$ and Picard number $>1$ are classified, see \cite[Cor.~3.1.15, Th.~3.3.1, Th.~7.2.15]{fanoEMS}. An inspection of the classification shows that the only cases where $\rho_Y=2$ and $Y$ has  two elementary contractions onto $\pr^2$ are $\pr^2\times\pr^2$ and 
a double cover of $\pr^2\times\pr^2$ ramified over a divisor of degree $(2,2)$.

On the other hand $Y$ cannot be covered by a family of curves of anticanonical degree $3$ by \cite[Lemma 3.12]{blowup}, and this excludes $\pr^2\times\pr^2$. In the other case
we have $h^0(Y,-K_Y)=45$ (see \href{https://arxiv.org/abs/2508.21207v1}{8.34}), and by \eqref{sections} this implies that $r\leq 2$ and $\rho_X\leq 4$, a contradiction.
\end{prg}
  \begin{prg}\label{duo}
We suppose from now on that $Y$ has index one.

For $i=1,2$ let us consider a $K$-negative resolution 
 $f_i\colon Y_i\to B_i$ of $Y\dasharrow B_i$;
we have a diagram:
$$\xymatrix{{Y_1}\ar[d]_{f_1}&Y\ar@{-->}[l]_{\psi_1}\ar@{-->}[r]^{\psi_2}&{Y_2}\ar[d]^{f_2}\\
{B_1}&&{B_2}
}$$
where $\psi_i$ is a SQM.
Fix $i\in\{1,2\}$ and let $D_i\subset Y_i$ be a general fiber of $f_i$ if $B_i=\pr^1$, the pullback of a general line if $B_i=\pr^2$.
Let $\w{D}_i\subset Y$ be the transform of $D_i$.
\end{prg}
  \begin{prg}
    We notice that
    the general fiber $F_i$ of $f_i$ cannot have a covering family of rational curves of anticanonical degree $3$.
    Otherwise, $Y_i$ would have a covering family $V$ of rational curves, of anticanonical degree $3$, contracted to points by $f_i$. However this is excluded as in \cite[4.24]{blowup}.

    This implies that $D_i$ contains an irreducible curve $\Gamma_i$, contracted by $f_i$,  with $c_i:=-K_{D_i}\cdot \Gamma_i\in\{1,2,4\}$ \cite[Cor.~IV.1.15]{kollar}, and $c_i\in\{1,2\}$ if $B_i=\pr^2$.  Notice that $D_i\cdot\Gamma_i=0$ and $-K_{Y_i}\cdot\Gamma_i=c_i$.

    Since $F_i$ is Fano, it is rationally connected.
Consider $f_i\colon Y_i\to B_i$ if $B_i=\pr^1$, or
$f_{i|D_i}\colon D_i\to\pr^1$ 
 if $B_i=\pr^2$. These are families of rationally connected varieties over a curve, and 
by \cite{GraberHarrisStarr} they have a section, namely
there exists $C_i\subset Y_i$ such that  $D_i\cdot C_i=1$.
    \end{prg}
  \begin{prg}\label{colorado}
We have $\Eff(Y)=\langle[\w{D}_1],[\w{D}_2]\rangle$, so there exist positive integers $n,a_1,a_2$ such that
$$n(-K_Y)=a_1\w{D}_1+a_2\w{D}_2,$$
and we can assume that $\gcd(n,a_1,a_2)=1$. 
Now let $\widehat{D}_2\subset Y_1$ be the transform of $D_2$. In $Y_1$ we have 
\stepcounter{thm}
\begin{equation}\label{Y1}
n(-K_{Y_1})=a_1D_1+a_2\widehat{D}_2,
\end{equation}
 and intersecting with $C_1$ we get 
$$n(-K_{Y_1}\cdot C_1)=a_1+a_2\widehat{D}_2\cdot C_1.$$
This implies that $\gcd(a_2,n)$ divides $a_1$, and hence that $\gcd(a_2,n)=1$. Similarly we get that $\gcd(a_1,n)=1$.

Set $d:=\gcd(a_1,a_2)$ and  $a_i:=da_i'$ for $i=1,2$. 
We have
$$n(-K_Y)=d(a'_1\w{D}_1+a'_2\w{D}_2)$$
and $\gcd(n,d)=1$. This implies that $d=1$, because $Y$ has index $1$. 
Now intersecting with $\Gamma_1$ in $Y_1$ we get from \eqref{Y1}:
$$nc_1=a_2\widehat{D}_2\cdot\Gamma_1,$$
which implies that $a_2$ divides $c_1$ and hence that $a_2\in\{1,2,4\}$, and $a_2\in\{1,2\}$ if $B_1=\pr^2$.
Similarly we deduce that $a_1\in\{1,2,4\}$.

Since $\gcd(a_1,a_2)=1$, 
up to exchanging $Y_1$ and $Y_2$ we can assume that $a_1=1$, so that \eqref{Y1} becomes:
\stepcounter{thm}
\begin{equation}\label{Y2}
  n(-K_{Y_1})=D_1+a_2\widehat{D}_2.
\end{equation}  
\end{prg}
  \begin{prg}
Suppose that $B_1\cong\pr^1$. Then $D_1$ is a smooth Fano $3$-fold and $h^0(D_1,-K_{Y_1|D_1})=h^0(D_1,-K_{D_1})=\frac{1}{2}(-K_{D_1})^3+3\leq 35$ \cite[Cor.~7.1.2]{fanoEMS}. Moreover by \eqref{Y2}:
$$-K_{Y_1}-D_1=a_2\widehat{D}_2-(n-1)(-K_{Y_1}),$$
and we know that $h^0(Y_1,-K_{Y_1})=h^0(Y,-K_Y)>0$ by Th.~\href{https://arxiv.org/abs/2508.21207v1}{2.12};
thus we get
\begin{gather*}
h^0(Y_1,-K_{Y_1}-D_1)=h^0\bigl(Y_1,a_2\widehat{D}_2-(n-1)(-K_{Y_1})\bigr)\leq 
h^0(Y_1,a_2\widehat{D}_2)
=h^0(Y_2,a_2D_2)\\=h^0\bigl(Y_2,f_2^*\ma{O}_{B_2}(a_2)\bigr)=h^0\bigl(B_2,\ma{O}_{B_2}(a_2)\bigr)\leq
h^0\bigl(\pr^2,\ma{O}_{\pr^2}(4)\bigr)=15.\end{gather*}

Now the exact sequence
$$0\la\ma{O}_{Y_1}(-K_{Y_1}-D_1)\la\ma{O}_{Y_1}(-K_{Y_1})\la\ma{O}_{D_1}(-K_{Y_1|D_1})\la 0$$
yields 
$$h^0(Y,-K_Y)=h^0(Y_1,-K_{Y_1})\leq h^0(Y_1,-K_{Y_1}-D_1)+h^0(D_1,-K_{Y_1|D_1})\leq 15+35=50,$$
that together with \eqref{sections} implies that $r\leq 3$ and $\rho_X\leq 5$, a contradiction.
\end{prg}
  \begin{prg}\label{verylast}
    Suppose now that $B_1\cong\pr^2$, and
consider the blow-up
of the first point
$\Bl_{p_1}Y\to Y$ in $\ph'\colon X'\to Y$. By \cite[Lemma 4.19(1)]{blowup}, applied to the map $X\dasharrow \Bl_{p_1}Y$, we see that there is a smooth Fano $4$-fold $W$, with $\rho_{W}=3$, and a SQM $\Bl_{p_1}Y\dasharrow W$. Moreover there is a SQM of $X$ which is obtained by blowing-up $W$ at $\rho_X-3$ distinct points, so that using again [\emph{ibid.}, Prop.~3.3]
we have:
 \stepcounter{thm}
 \begin{equation}\label{sections2}h^0(W,-K_{W})=h^0(X,-K_X)+15(\rho_X-3).\end{equation}
 
 Since $p_1$ cannot belong to an exceptional plane [\emph{ibid.}, Lemma 4.4(3)], we see that $p_1\in\dom(\psi_i)$ for $i=1,2$ (see \ref{duo}); we still denote by $p_1$ its image in $Y_i$. Let
 $$\sigma_i\colon W_i\la Y_i$$ be the blow-up of $p_1$, and $E_i\subset W_i$ the exceptional divisor.
Then $W_i$ is another SQM of $\Bl_{p_1}Y$ and $W$.

Consider the composition $f_1\circ\sigma_1\colon W_1\to\pr^2$. As in \ref{avocado}, we see that this map must factor through an elementary rational contraction $W_1\dasharrow \mathbb{F}_1$, where $h\colon \mathbb{F}_1\to \pr^2$ is the blow-up of the point $f_1(p_1)$. 
$$\xymatrix{ {W_1}\ar@{-->}[d]\ar[r]^{\sigma_1}&{Y_1}\ar[d]^{f_1}\\
 {\mathbb{F}_1}\ar[r]^h&{\pr^2}
  }$$
  Recall that $D_1=f_1^*(l)$ where  $l\subset\pr^2$ is a general line (see \ref{duo}).  Using \eqref{Y2} we get:
  \stepcounter{thm}
  \begin{equation}\label{summer}
  \begin{split}
  -nK_{W_1} = & \ \sigma_1^*(-nK_{Y_1})-3nE_1=\sigma_1^*\bigl(D_1+a_2\wi{D}_2\bigr)-3nE_1\\
 = &\ \sigma_1^*\bigl(f_1^*(l)\bigr)+a_2\sigma_1^*(\wi{D}_2)-3nE_1.
  \end{split}
 \end{equation}

  Let $\alpha\colon W_1'\to\mathbb{F}_1$ be a $K$-negative resolution of $W_1\dasharrow\mathbb{F}_1$, $\wi{D}_2'\subset W_1'$ the transform of $\sigma_1^*(\wi{D}_2)$, and
$E_1'\subset W_1'$ the transform of $E_1=\Exc(\sigma_1)$. 
We have
  $E_1'=\alpha^*(e)$ where  $e\subset \mathbb{F}_1$
  is the $(-1)$-curve. Let also
  $\pi\colon\mathbb{F}_1\to\pr^1$ be the $\pr^1$-bundle; we have $h^*(l)=f+e$ where $f\subset \mathbb{F}_1$ is a general fiber of $\pi$. Finally let us consider $F=\alpha^*(f)$  a general fiber of $\pi\circ\alpha\colon W_1'\to\pr^1$.
$$\xymatrix{{W_1'}\ar@{-->}[r]\ar[dr]^{\alpha}& {W_1}\ar@{-->}[d]\ar[r]^{\sigma_1}&{Y_1}\ar[d]^{f_1}\\
{\pr^1}& {\mathbb{F}_1}\ar[l]_{\pi}\ar[r]^h&{\pr^2}
}$$

We note that $F$ is a smooth Fano $3$-fold. Indeed there is a $K$-negative resolution $W_1''\to\pr^1$ of $\pi\circ\alpha$, and the indeterminacy locus of the SQM $W_1'\dasharrow W_1''$ is contracted to points by $\alpha$, therefore it does not meet $F$. Hence $F$ is isomorphic to a general fiber of  $W_1''\to\pr^1$, which is Fano. We also note that $\alpha(F)=f$, so that $\rho_F\geq 2$.
  
In $W_1'$ we get, from \eqref{summer}: 
 \begin{align*}  -nK_{W_1'}=& \ \alpha^*\bigl(h^*(l)\bigr)+a_2\wi{D}_2'-3nE_1'=
                              \alpha^*(f+e)+a_2\wi{D}_2'-3nE_1'\\
   = & \  F+a_2\wi{D}_2'-(3n-1)E_1'.
 \end{align*}
 Recall that $\ol_{Y_2}(D_2)=f_2^*\ol_{B_2}(1)$ (see \ref{duo}), and
 that $a_2\in\{1,2\}$ because $B_1\cong\pr^2$ (see \ref{colorado}). 
  Therefore we have:
  \begin{gather*}
    h^0(W_1',-K_{W_1'}-F)= h^0\bigl(W_1',a_2\wi{D}_2'-(3n-1)E_1'-(n-1)(-K_{W_1'})\bigr)\\ \leq  h^0(W_1',a_2\wi{D}_2')
    =h^0\bigl(W_2,\sigma_2^*(a_2D_2)\bigr)
  =h^0\bigl(W_2,\sigma_2^*\bigl(f_2^*\ol_{B_2}(a_2)\bigr)\bigr)\\ =h^0(B_2,\ol_{B_2}(a_2))\leq  h^0(\pr^2,\ol(2))=6.
\end{gather*}
Moreover $-K_{W_1'|F}=-K_F$ and $h^0(F,-K_F)=\frac{1}{2}(-K_F)^3+3\leq 34$ by classification \cite[Ch.~12]{fanoEMS}, and finally we get $h^0(W,-K_W)=h^0(W_1',-K_{W_1'})\leq  h^0(W_1',-K_{W_1'}-F)+h^0(F,-K_F)\leq 40$. Together with \eqref{sections2}
this implies that $\rho_X\leq 5$, again a contradiction. This concludes the proof of Prop.~\ref{all30}.\qedhere
\end{prg}
\end{proof}

\begin{thebibliography}{99}
\bibitem[Cas17]{blowup} \bysame, \emph{{F}ano 4-folds, flips, and blow-ups of points}, J.~Algebra
  \textbf{483} (2017), 362--414.
\bibitem[IP99]{fanoEMS} V.A. Iskovskikh and Yu.G. Prokhorov, \emph{Algebraic geometry {V} - {F}ano
  varieties}, Encyclopaedia Math.\ Sci.\, vol.~47, Springer-Verlag, 1999.
\bibitem[Kol96]{kollar} J.~Koll{\'a}r, \emph{Rational curves on algebraic varieties}, Ergebnisse der
  Mathematik und ihrer Grenzgebiete, vol.~32, Springer-Verlag, 1996.
\bibitem[GHS03]{GraberHarrisStarr} T.~Graber, J.~Harris, and J.~Starr, \emph{Families of rationally connected
  varieties}, J.\ Amer.\ Math.\ Soc.\ \textbf{16} (2003), 57--67.
\end{thebibliography}
\end{document}
